% 本文件由 LaTeX 排版 Agent 根据有效 translation.json 自动生成。
% author=Councils; work_id=3809; title=君士坦丁堡会议（公元382年）

\chapter{君士坦丁堡会议（公元382年）}

% structure: p0001 source=h1 target=omitted reason=page-title-already-rendered-by-parent

% structure: p0002 source=h2 target=section reason=relative-heading-rank; base=h2

\section{会议信函}

致最可敬的诸位上主、我们的极可敬的诸位弟兄及同僚\Person{达玛稣}、\Person{盎博罗削}、\Person{布里东}、\Person{华肋廉}、\Person{阿所留}、\Person{阿乃缪}、\Person{巴西略}，以及聚集于\Place{罗马}大城的其余诸位圣主教：聚集于\Place{君士坦丁堡}大城的正统主教神圣会议，在主内致候。\par

若要逐一陈述我们因亚略派势力所受的一切苦难，并试图向你们诸位尊驾提供消息，仿佛你们尚不熟知，则似为多余。因为我们并不以为你们的虔诚将我们所遭遇之事视为如此次要，以致你们需要得知那些必然唤起同情之事。况且，袭击我们的风暴也并非微小到可以避人耳目。我们的迫害不过是昨日之事。那声音仍在遭受迫害者以及那些因爱德而以他人痛苦为己苦者的耳中回响。可以说，就在一两天前，有些从异乡解除锁链的人，历经重重苦难返回了自己的教会；另有一些死于流放者的遗骸被运回故乡；还有些人，即使从流放地返回后，仍见异端者的热情沸腾，竟如真福\Person{斯德望}一般被他们用石头砸死，在自己的土地上遭遇了比异乡更悲惨的命运。还有一些人，因各种残忍手段而憔悴，仍在其身上带着伤疤和基督的烙印。谁能述说罚款、褫夺公权、个人没收财产、阴谋、凌辱、监禁的故事？实在，各种磨难在我们身上无计其数地发生，也许是我们在为罪过付上代价，也许是仁慈的天主藉着我们多重的苦难来试炼我们。为此，我们一切感谢天主，祂藉着这样的磨难训练了祂的仆人，并按照祂丰厚的仁慈再次使我们得到安息。我们确实需要长久闲暇、时间与辛劳来重建教会，好使她能像医生长期治愈身体、逐步驱除疾病一样，复归真正宗教的古老健康。诚然，大体上我们似乎已从迫害的暴力中得救，如今正在收复那长期为异端者所掠夺的教会。但仍有豺狼困扰我们，它们虽被逐出羊栈，却在林间蹂躏羊群，胆敢举行对立集会，在百姓中煽动叛乱，且无所不为以损害教会。因此，正如我们已说过的，我们必须更长久地劳苦。然而，既然你们曾以兄弟之爱邀请我们（如同我们是你们自己的肢体），藉我们最虔诚的皇帝的书信，前往你们蒙天主恩准在\Place{罗马}召开的会议，目的是：既然我们当时独受迫害之苦，如今当我们的皇帝在真正宗教上与我们同心时，你们不可离我们而统治，而应使我们，用宗徒的话说，与你们一同为王。我们的祈祷是，若可能，全体一同离开我们的教会，宁可满足我们渴望见你们的心愿，也不顾及教会的需要。因为谁给我们鸽子般的翅膀？我们便要飞去安息。但此途似乎会使刚刚恢复的教会毫无防卫，而且对多数人而言亦不可行，因为按照一年前你们从\Place{阿奎莱亚}会议后寄给最虔诚的皇帝\Person{德奥多西}的书信，我们已前往\Place{君士坦丁堡}，只备足了抵达\Place{君士坦丁堡}的旅费，并携带着仅此会议各省留守主教的同意。我们未曾期待更远的旅程，在抵达\Place{君士坦丁堡}之前亦未闻一词。除此以外，因指定时间之短促，既无法预备更长旅程，也无法与我们共融的各省主教联络并获得他们的同意，因此对大多数人而言，前往\Place{罗马}的旅程是不可能的。因此，我们采取了在当前情况下最可行的办法，既为更好地管理教会，也为显明我们对你们的爱德，极力敦促我们最可敬、最可贵的同僚及弟兄主教\Person{西里亚谷}、\Person{欧瑟比}和\Person{普利西安}，同意前往你们那里。\par

藉着他们，我们愿表明：我们的心意全然是和平，以合一为唯一目标，且对正确的信德充满热忱。因为我们无论遭受迫害、磨难、帝王的威胁、王公的残暴，或任何来自异端者的试炼，皆为福音的信德忍受一切；这信德由在\Place{彼提尼雅}的\Place{尼西亚}的三百一十八位教父所批准。这信德应当为你们、为我们、为一切不曲解真正信德之言的人所满足；因为它是古老的信德，是我们圣洗的信德，是教导我们相信圣父、圣子及圣神之名的信德。按照这信德，圣父、圣子及圣神具有同一天主性、能力和本体；尊荣相等，威严相等，存在于三个完美的位格，即三个完美的位格。如此，\Person{撒伯里乌}的异端因混淆位格（即毁灭位格性）而无立足之地；\Person{欧诺米}派、\Person{亚略}派及反圣神派的亵渎亦被废止，它们分裂本体、本性及天主性，并在非受造的、同体的、同永的天主圣三之上引入一种后生的、受造的、不同本体的本性。此外，我们保守主的降生成人的道理不受歪曲，持守肉身救恩计划既非无灵魂、无理智，亦非不完全的传授；且深知天主的圣言在万世之前已是完全的，在末世为了我们的救恩成为完全的人。\par

关于我们无畏且坦率宣讲的教义概要，就此足矣。若蒙您垂询，亦可查阅安提约基雅主教会议之卷册，以及去年在君士坦丁堡举行之大公会议所颁卷册，其中我们更详尽地陈明了信仰宣认，并附加了绝罚，针对近来创新者所记之异端。\par

至于个别教会的具体管理，如您所知，一项古老习俗业已确立，并经尼西亚圣教父之谕令确认：各教省主教，连同其同意之邻省主教，应按实际需要主持祝圣。依循此等习俗，其他教会均由我们及最著名教会的司铎公开委任管理。因此，对于新建（若可如此说）的君士坦丁堡教会，我们最近藉天主仁慈从异端者的亵渎中救出（犹如从狮口夺出），在大公会议中，在至虔敬的德奥多西皇帝面前，经全体圣职人员和全城同意，共同祝圣了可敬且至虔诚的奈克塔留斯为主教。至于叙利亚那最古老且真正宗徒的教会——那里首先被赋予基督徒这崇高尊号——教省主教及东部教区的主教会集，依法祝圣了可敬且至虔诚的弗拉维雅努斯为主教，得全体教会同意，他们如同一心表达对他的尊敬。此合法祝圣也得到大公会议的认可。至于耶路撒冷教会——诸教会之母——我们宣告可敬且至虔诚的济利禄为主教，他先前已由教省主教依法祝圣，并多次与亚略异端者英勇奋战。我们恳请尊座为此等由我们依法依规所定之事而喜乐，藉着属灵之爱与敬畏上主之情，约束人情，使教会之建立重于个人恩宠或情面。如此，因我们中间有信仰的一致并建立了基督徒的爱德，我们将不再使用宗徒所谴责的说法：\BibleQuote{我是属保禄的，我是属阿颇罗的，我是属刻法的}，而全体显为属基督的，他在我们中间并未分裂，我们将藉天主的恩宠保持教会身体不分裂，并勇敢地站立于主的审判座前。\par

\SourceRecord{Translated by Henry Percival.；Nicene and Post-Nicene Fathers, Second Series；Vol. 14.；Edited by Philip Schaff and Henry Wace.；Buffalo, NY: Christian Literature Publishing Co.,；1900.；<http://www.newadvent.org/fathers/3809.htm>.}
